\documentclass[11pt, a4paper]{article}

% Gemeinsame Style-Definitionen (TEXINPUTS muss templates/ enthalten — siehe scripts/build_tex.py)
\input{bfw-style.tex}

\title{\vspace*{-1.5cm}\Huge\bfseries\color{bfwPrimaryDark}Session~4: Rechtsformen \& Aufbauorganisation\\[0.4em]
  \large\color{bfwInkSoft}\textnormal{Von Einzelunternehmen bis AG --- Leitungssysteme und Vollmachten im Überblick}}
\author{\textsf{IT Lernfeld 1 --- Das Unternehmen und die eigene Rolle im Betrieb beschreiben}}
\date{}

\begin{document}
\maketitle
\thispagestyle{empty}

\vspace{1em}
\begin{merksatz}[title={\bfseries Worum geht es in dieser Session?}]
Jedes Unternehmen muss sich für eine \emph{Rechtsform} entscheiden --- sie bestimmt,
wer haftet, wer leitet, wie viel Kapital benötigt wird und wie das Unternehmen besteuert
wird. Am Beispiel der \emph{JIKU IT-Solutions GmbH} wird deutlich, warum die GmbH die
häufigste Wahl bei IT-Systemhäusern ist. Daneben klärt diese Session, was eine
\emph{Firma} im juristischen Sinn ist, welche Vollmachten (\emph{Prokura} und
\emph{Handlungsvollmacht}) es gibt und wie sich ein Unternehmen durch
\emph{Aufbauorganisation} gliedert.
\end{merksatz}

\tableofcontents
\vspace{1em}
\hrule
\vspace{1em}

\section{Aufbau- und Ablauforganisation}

\subsection{Bereiche betrieblicher Organisation}

\begin{center}
\begin{tabular}{p{6cm} p{7cm}}
\toprule
\textbf{Aufbauorganisation} & \textbf{Ablauforganisation}\\
\midrule
Gliederung des Betriebes in funktionsfähige Einheiten; Übertragung von Aufgaben, Kompetenzen
und Verantwortung; Festlegung von Informations- und Anweisungswegen &
Regelt den funktionalen, zeitlichen und räumlichen Arbeitsablauf zum optimalen Einsatz
vorhandener Kapazitäten an Mitarbeitern und Betriebsmitteln\\
\bottomrule
\end{tabular}
\end{center}

\subsection{Stellenarten}

\begin{center}
\begin{tabular}{p{3.5cm} p{9.5cm}}
\toprule
\textbf{Stellenart} & \textbf{Definition}\\
\midrule
Instanz & Stelle mit Anordnungs- und Entscheidungsbefugnis gegenüber rangniedrigeren Stellen\\
\addlinespace
Stabsstelle & Hilfsstelle einer Instanz \emph{ohne} Weisungsbefugnis; beratungs- und
entscheidungsvorbereitende Funktion; hat Informationsrecht aus allen Linienstellen;
Beispiele: Assistent der Geschäftsleitung, Qualitätsmanager, Datenschutzbeauftragter\\
\addlinespace
Ausführende Stelle & Ohne Leitungs- und Entscheidungsbefugnis\\
\bottomrule
\end{tabular}
\end{center}

\begin{merksatz}[title={\bfseries Merke: Stabsstelle}]
Eine Stabsstelle \textbf{berät} die Instanz und bereitet Entscheidungen vor --- sie darf jedoch
\textbf{keine Weisungen} erteilen. JIKU IT-Solutions hat laut Organigramm mehrere Stabsstellen
(z.\,B.\ IT-Sicherheit, Qualitätsmanagement).
\end{merksatz}

\section{Leitungssysteme und Organigramm}

Das \textbf{Organigramm} ist ein Schaubild des Leitungssystems; es zeigt Hierarchieebenen,
Über-/Unterordnungsbeziehungen und Stabsstellen.

\begin{center}
\begin{tabular}{p{3.5cm} p{4.5cm} p{4.5cm}}
\toprule
\textbf{System} & \textbf{Stärken} & \textbf{Schwächen}\\
\midrule
Einliniensystem & Klare Zuständigkeiten; überschaubare Struktur & Lange Weisungswege; Überlastung höherer Instanzen; mangelnde Flexibilität\\
\addlinespace
Stabliniensystem & Vorteile der Linie + bessere Beratung durch Stabsstellen & Kompetenzkonflikte zwischen Stabsstellen und Linieninstanzen\\
\addlinespace
Mehrliniensystem & Flexibler Mitarbeitereinsatz; Spezialisierung möglich & Abgrenzungsprobleme; Mitarbeiter können überfordert werden\\
\addlinespace
Matrixorganisation & Ergebnis durch zwei Kernkompetenzen verbessert; Teamarbeit gefördert & Zuständigkeitskonflikte; ungleiche Arbeitsbelastung möglich\\
\bottomrule
\end{tabular}
\end{center}

\begin{merksatz}[title={\bfseries JIKU IT-Solutions: Kombination aus Stablinien- und Matrixorganisation}]
Die Geschäftsleitung von JIKU hat sich für eine Kombination aus \textbf{Stabliniensystem}
und \textbf{Matrixorganisation} entschieden. In den kundennahen IT-Solutions-Bereichen wird
viel im Projektmanagement gearbeitet, sodass dort eine Matrixorganisation am besten passt.
\end{merksatz}

\subsection{Funktionale und divisionale Organisation}

\begin{center}
\begin{tabular}{p{3.5cm} p{9.5cm}}
\toprule
\textbf{Form} & \textbf{Beschreibung}\\
\midrule
Funktionale Organisation & Zweite Hierarchiestufe nach betrieblichen Funktionen (z.\,B.\ Beschaffung,
Fertigung, Absatz, Verwaltung); Koordination und Steuerung nach Abteilungen\\
\addlinespace
Spartenorganisation (divisional) & Zweite Hierarchiestufe nach Produktbereichen oder Sparten mit
eigener Kompetenz; Profit- oder Costcenter möglich\\
\bottomrule
\end{tabular}
\end{center}

\section{Kaufmannseigenschaft und Handelsregister}

Wer ein Handelsgewerbe betreibt oder in das \textbf{Handelsregister} eingetragen ist, erlangt
die \textbf{Kaufmannseigenschaft} und unterliegt neben dem BGB auch dem HGB (doppelte Buchführung,
Firmenpflicht, Prokura u.\,a.).

\begin{center}
\begin{tabular}{p{3cm} p{10cm}}
\toprule
\textbf{Merkmal} & \textbf{Inhalt}\\
\midrule
Zweck & Offenlegung wesentlicher Rechtsverhältnisse; öffentliche Informationsquelle für die Wirtschaft\\
Abteilung A & Einzelkaufleute und Personengesellschaften (OHG, KG)\\
Abteilung B & Kapitalgesellschaften (GmbH, AG)\\
Schutzfunktion & Schutz von Kunden und Lieferanten durch Bekanntmachung der Rechtsverhältnisse\\
Publikationsfunktion & Jedermann kann Einsicht nehmen; Veröffentlichung auch im Bundesanzeiger\\
\bottomrule
\end{tabular}
\end{center}

\section{Firma und Firmengrundsätze (§§ 17 ff. HGB)}

Die \textbf{Firma} ist nach \textbf{§~17 HGB} der Name, unter dem ein Kaufmann seine Geschäfte
betreibt und die Unterschrift abgibt. Nur eingetragene Kaufleute dürfen eine Firma führen;
Kleingewerbetreibende verwenden lediglich eine \emph{Geschäftsbezeichnung}.

\begin{center}
\begin{tabular}{p{3cm} p{10cm}}
\toprule
\textbf{Firmenart} & \textbf{Beschreibung und Beispiel}\\
\midrule
Personenfirma & Familienname des Inhabers, z.\,B.\ Lasse Müller e.\,Kfm.\\
Sachfirma & Unternehmensgegenstand, z.\,B.\ IT-Sander e.\,K., Sporti~GmbH\\
Fantasiefirma & Freie Namensgestaltung, z.\,B.\ Sysygy e.\,Kfm., IT-2000 GmbH\\
\bottomrule
\end{tabular}
\end{center}

\begin{center}
\begin{tabular}{p{3cm} p{10cm}}
\toprule
\textbf{Grundsatz} & \textbf{Inhalt}\\
\midrule
Firmenausschließlichkeit & Die Firma muss sich am selben Ort von anderen Firmen unterscheiden (§~30 HGB).\\
Firmenwahrheit & Irreführungsverbot: keine täuschenden Angaben über die geschäftlichen Verhältnisse (§~18 Abs.~2 HGB).\\
Firmenbeständigkeit & Die Firma kann auch bei Inhaberwechsel fortgeführt werden; Schutz durch §~12 BGB, §~37 HGB, §~15 MarkenG.\\
\bottomrule
\end{tabular}
\end{center}

\section{Prokura und Handlungsvollmachten}

Nur im Handelsregister eingetragene Unternehmen können Prokura und Handlungsvollmachten
nach dem HGB vergeben.

\begin{center}
\begin{tabular}{p{3.5cm} p{5.5cm} p{4cm}}
\toprule
\textbf{Merkmal} & \textbf{Prokura (§§ 48–53 HGB)} & \textbf{Handlungsvollmacht (§§ 54–58 HGB)}\\
\midrule
Umfang & Alle gewöhnlichen und außergewöhnlichen Geschäfte; Ausnahmen: Bilanz/Steuererklärung,
HR-Eintragungen, Prokura erteilen/entziehen, Gesellschafter aufnehmen &
Gewöhnliche Handelsgeschäfte (Gesamtvollmacht); bestimmte Art (Artvollmacht); Einzelgeschäft
(Einzelvollmacht)\\
\addlinespace
Zeichnung & \texttt{ppa.~Name} oder \texttt{p.p.a.~Name} & \texttt{i.V.~Name} (Gesamt-/Artvollmacht) bzw.\ \texttt{i.A.~Name} (Einzelvollmacht) \newline \textit{Hinweis: i.V. und i.A. sind kaufmännische Konventionen, gesetzlich nicht vorgeschrieben.}\\
\addlinespace
Erteilung & Nur durch Kaufmann (Inhaber/Geschäftsführer); Eintragung ins HR vorgeschrieben
(deklaratorisch) & Formlos (mündlich, schriftlich); keine HR-Eintragung möglich\\
\addlinespace
Widerruf & Jederzeit möglich & Jederzeit möglich\\
\bottomrule
\end{tabular}
\end{center}

\section{Rechtsformen im Überblick}

Entscheidungskriterien für die Rechtsformwahl sind u.\,a.: Haftung, Leitungsbefugnis,
Gewinn-/Verlustbeteiligung, Kapitalbedarf, Flexibilität bei Gesellschafterwechsel,
Steuerbelastung, Offenlegungspflichten und Gründungsaufwand.

\begin{center}
\begin{tabular}{p{2.2cm} p{2.5cm} p{2.5cm} p{2.5cm} p{3cm}}
\toprule
\textbf{Merkmal} & \textbf{Einzelunternehmen} & \textbf{GbR} & \textbf{OHG} & \textbf{KG}\\
\midrule
Gründer & 1 Inhaber & mind.\ 2 & mind.\ 2 & mind.\ 1 Komplementär + 1 Kommanditist\\
Mindestkapital & keine Untergrenze & keine Untergrenze & keine Untergrenze & keine Untergrenze\\
Haftung & persönlich, unbeschränkt & persönlich, unbeschränkt, solidarisch & persönlich, unbeschränkt, solidarisch & Komplementär: persönlich, unbeschränkt; Kommanditist: bis zur Einlage\\
Leitung & Inhaber allein & alle Gesellschafter & alle Gesellschafter & nur Komplementär\\
Rechtsgrundlage & § 14 BGB, § 19 HGB & §§ 705 ff.\ BGB & §§ 105 ff.\ HGB & §§ 161 ff.\ HGB\\
\bottomrule
\end{tabular}
\end{center}

\begin{center}
\begin{tabular}{p{2.2cm} p{3.5cm} p{4cm} p{3cm}}
\toprule
\textbf{Merkmal} & \textbf{GmbH} & \textbf{UG (haftungsbeschränkt)} & \textbf{AG}\\
\midrule
Gründer & mind.\ 1 Gesellschafter & mind.\ 1 Gesellschafter & mind.\ 1 Aktionär\\
Mindestkapital & \textbf{25.000~€} Stammkapital (§~5 GmbHG); 50\,\% sofort einzuzahlen & \textbf{ab 1~€} Stammkapital; 25\,\% Thesaurierungspflicht & \textbf{50.000~€} Grundkapital (§~7 AktG)\\
Haftung & nur die Gesellschaft haftet & nur die Gesellschaft haftet & nur die Gesellschaft haftet\\
Leitung & Geschäftsführer (notariell beurkundet); Gesellschafterversammlung & Geschäftsführer & Vorstand, Aufsichtsrat, Hauptversammlung\\
Rechtsgrundlage & GmbH-Gesetz & GmbH-Gesetz & Aktiengesetz\\
\bottomrule
\end{tabular}
\end{center}

\begin{merksatz}[title={\bfseries Merke: Mindestkapital der Kapitalgesellschaften}]
UG (haftungsbeschränkt): \textbf{ab 1~€} Stammkapital --- 25\,\% des Jahresüberschusses
müssen thesauriert werden, bis 25.000~€ erreicht sind.\\[0.3em]
GmbH: \textbf{25.000~€} Mindest-Stammkapital (§~5 GmbHG); bei Gründung mindestens
12.500~€ sofort einzuzahlen.\\[0.3em]
AG: \textbf{50.000~€} Mindest-Grundkapital (§~7 AktG).
\end{merksatz}

\begin{merksatz}[title={\bfseries Steuerliche Einordnung}]
Personengesellschaften (Einzelunternehmen, GbR, OHG, KG): Jeder Gesellschafter zahlt
\textbf{Einkommensteuer} auf seinen Gewinnanteil (§§ 15, 18 EStG).\\[0.3em]
Kapitalgesellschaften (GmbH, AG): Die Gesellschaft ist als juristische Person eigenes
Steuersubjekt und zahlt \textbf{Körperschaftsteuer} auf ihren Gewinn (§ 1 KStG).
\end{merksatz}

\end{document}
